\section{Methodology}
\label{sec:methodology}
\subsection{Data and Scope}
\label{sec:data}
We analyze the released productivity tables associated with
Guglielmi et al.~\cite{guglielmi2026copilot}. They contain 16 experts and
17 novices, each with one session per mode on different Java maintenance tasks.
We retain the source experience labels and all 66 released observations. The
source investigators collected the data; we recruit no participants and collect
no new human outcomes. The replication-package link, attribution, and SHA-256
checksums are recorded in the repository's data provenance files.

Figure~\ref{fig:pipeline} summarizes our analytical pipeline. Table~\ref{tab:parameters}
lists its implemented parameters. We verify file hashes, normalize decimal-comma
numbers and time-column spellings, require finite values and positive times,
and check completion bounds. We compare recalculated productivity with the
published value using an absolute tolerance of $10^{-8}$. Each participant must
have exactly one observation per mode, consistent group metadata, and different
tasks. Invalid inputs stop processing; they are not silently excluded.

\begin{figure}[t]
\centering
\begingroup\small
\fbox{\begin{minipage}{0.91\columnwidth}
\centering
Published CSV tables and provenance\\
$\downarrow$\\
Checksum, schema, and participant-pair validation\\
$\downarrow$\\
Productivity recalculation and paired changes\\
$\downarrow$\\
Group summaries and stratified bootstrap\\
$\downarrow$\\
CSV summaries and JSON results
\end{minipage}}
\endgroup
\caption{Implemented secondary-analysis pipeline. Metric calculation occurs
after validation. These are participant outcomes, not runtime benchmark logs.}
\label{fig:pipeline}
\end{figure}

\begin{table}[t]
\caption{Implemented analysis parameters in \texttt{configs/config.json}.}
\label{tab:parameters}
\centering
\begin{tabular}{@{}ll@{}}
\toprule
Parameter & Value \\
\midrule
Participants & 33 (17 novice, 16 expert) \\
Observations & 66 sessions \\
Bootstrap repetitions & 10,000 \\
Pseudorandom seed & 42 \\
Confidence level & 95\% \\
Completion-drop flag & Greater than 5 percentage points \\
\bottomrule
\end{tabular}
\end{table}

\subsection{Outcomes and Estimands}
\label{sec:estimands}
Equation~\ref{eq:productivity} defines productivity $P_{im}$ for participant $i$
under mode $m$, where $c_{im}$ is completion in percent and $t_{im}$ is time in
minutes:
\begin{equation}
P_{im}=\frac{c_{im}/100}{t_{im}/60}.
\label{eq:productivity}
\end{equation}
The unit is completed task fraction per hour. Incomplete and zero-completion
sessions remain in the analysis. Session time is not necessarily time to a fully
correct solution. Completion does not measure security or comprehensive code quality.

Equation~\ref{eq:paired} defines the automatic-minus-manual change, and
Equation~\ref{eq:contrast} defines the primary expert-minus-novice contrast:
\begin{equation}
D_i=P_{i,\mathrm{automatic}}-P_{i,\mathrm{manual}}.
\label{eq:paired}
\end{equation}
\begin{equation}
\theta=\bar{D}_{E}-\bar{D}_{N},\qquad
\bar{D}_{g}=\frac{1}{n_g}\sum_{i\in g}D_i.
\label{eq:contrast}
\end{equation}
Here $E$ and $N$ denote expert and novice groups and $n_g$ is their participant
count. Positive $\theta$ indicates a larger automatic-mode benefit for experts.
RQ1 and RQ2 concern $\bar{D}_{N}$ and $\bar{D}_{E}$, respectively; RQ3 concerns
$\theta$. All comparisons are two-sided and exploratory.

We additionally calculate paired completion and time changes. Our student-defined
completion guardrail flags a group mean change below $-5$ percentage points.
It is a reporting rule, not a validated quality threshold. No API error rate or
technical warm-up exclusion is inferred from human session measurements.

\subsection{Uncertainty and Reproducibility}
\label{sec:uncertainty}
The implemented analysis reports group--mode means, medians, and 95th
percentiles. For paired mean changes and $\theta$, it generates 10,000 bootstrap
samples using seed 42. Sampling occurs separately within experience groups and
retains each participant's pair. The interval endpoints are the 2.5th and 97.5th
percentiles of the resulting estimates, with linear interpolation. Sessions are
not resampled as independent participants.

The pipeline uses the Python standard library. The repository documents local
Python 3.12 and a Docker runtime pinned to Python 3.12.12. Running
\texttt{src/benchmark.py} with the full-data configuration generates validated
observations, paired changes, and JSON summaries. These are unadjusted estimates.
Task/order adjustment, leave-one-participant-out checks, the executed analysis
notebook, and final plots remain Task 3 work; they are not presented as completed
methods in this draft.
